\documentclass[11pt]{amsart}
\usepackage[margin=1.15in]{geometry}
\usepackage{amsmath,amssymb,amsthm,booktabs}
\usepackage[colorlinks=true,allcolors=blue]{hyperref}
\newtheorem{theorem}{Theorem}
\newtheorem{lemma}[theorem]{Lemma}
\newtheorem{proposition}[theorem]{Proposition}
\newtheorem{question}[theorem]{Question}
\theoremstyle{remark}
\newtheorem{remark}[theorem]{Remark}
\newcommand{\tr}{\operatorname{tr}}
\begin{document}
\title[Removing edges from flower snark line graphs]{How many edges can be removed from a flower snark line graph?}
\author{Zijian Zeng}
\address{UCSI University, Kuala Lumpur, Malaysia}
\date{17 November 2026}

\begin{abstract}
The line graph $L(J_q)$ of the flower snark $J_q$ has $6q$ vertices and $12q$ edges, and it is not
$3$-colourable when $q$ is odd. Exercise 7.2.2.3--233 of Knuth's \emph{The Art of Computer Programming}
asks how many of its edges can be removed before it becomes $3$-colourable; Knuth's best answer so far
removes nine. We show that nine is optimal for every odd $q\ge5$: no ten edges can be removed, and the
removable nine-edge sets are exactly the $4q$ sets obtained by isolating the three edges of $J_q$ at one vertex.
For $q=3$ the maximum is fifteen. The lower bound holds for every snark by the parity lemma. The upper bound
uses a $27$-state transfer matrix $T$ around the ring of $J_q$; since $T^2=T^4$, the problem for all $q$
reduces to a finite computation, which we carry out and check in several independent ways. We also find the
number of removable $k$-edge sets for every $k$: it is $q$ times $12,132,324,498,504,336,144,36,4$ for
$k=1,\dots,9$ when $q\ge7$.
\end{abstract}
\maketitle

\section{The question and the answer}
In this note a \emph{snark} is a connected simple cubic graph whose edges cannot be coloured with three
colours so that the three edges at every vertex get different colours (we do not require girth at least five,
so that $J_3$, which has a triangle, is included). The flower snarks $J_q$ of Isaacs~\cite{isaacs}, $q\ge3$
odd, have vertices $a_i,b_i,c_i,d_i$ for $i\in\mathbb Z_q$; each centre $a_i$ is joined to $b_i,c_i,d_i$; the
$b_i$ form a $q$-cycle $b_0b_1\ldots b_{q-1}b_0$; and the $c_i$ and $d_i$ together form one $2q$-cycle
$c_0c_1\ldots c_{q-1}d_0d_1\ldots d_{q-1}c_0$. Its line graph $L(J_q)$ has the $6q$ edges of $J_q$ as vertices,
two of them adjacent when they share an endpoint; so $L(J_q)$ is the union of $4q$ triangles, one for each
vertex of $J_q$, and has $12q$ edges. A proper $3$-colouring of $L(J_q)$ is the same thing as a $3$-edge-colouring
of $J_q$, so $L(J_q)$ is not $3$-colourable when $q$ is odd.

Knuth uses these graphs as test cases for SAT solvers~\cite[\S7.2.2.2, ex.~176]{knuth4b}, \cite[\S7.2.2.3]{knuth7}.
His exercise 7.2.2.3--233~\cite{knuth7}, rated~$40$ (numbering as in the draft cited), observes that $L(J_q)$ stays non-$3$-colourable when any one
edge is removed, and asks for the largest number of edges whose removal leaves it non-$3$-colourable.
His answer says that he hopes a reader will supply a good answer, and that his best so far is to remove nine
edges; his example removes the nine edges of $L(J_q)$ that touch the triangle $\{b_1,c_1,d_1\}$ of his
labelling. (His labels name the vertices of $L(J_q)$, i.e.\ the edges of $J_q$; judging from the adjacencies in
his example, $b_1,c_1,d_1$ are the three spokes at the centre of unit~$1$, so in our notation the example is the
claw set $R_{a_1}$ defined below.)

Call a set $R$ of edges of $L(J_q)$ \emph{removable} if $L(J_q)-R$ is not $3$-colourable. A subset of a removable
set is removable. For a vertex $v$ of $J_q$ let $R_v$ be the set of edges of $L(J_q)$ that touch at least one of
the three edges of $J_q$ at $v$: the three edges of the triangle of $v$ and the two further edges at each of its
corners, nine in all. We call $R_v$ the \emph{claw set} of $v$. Knuth's example is a claw set.

\begin{theorem}\label{thm:main}
Let $q\ge5$ be odd. Every removable set of edges of $L(J_q)$ has at most nine edges, and the removable sets with
nine edges are exactly the $4q$ claw sets $R_v$. For $q=3$ the largest removable set has fifteen edges, and it
is unique: it consists of the edges that touch the six edges of $J_3$ at $b_0$, $b_1$, $b_2$.
\end{theorem}

So the answer to the exercise is \emph{nine}, for every odd $q\ge5$, and Knuth's example is best possible. We
also count all removable sets.

\begin{theorem}\label{thm:count}
For odd $q\ge7$ the number of removable $k$-edge sets of $L(J_q)$ is $c_k\,q$, where
\[(c_1,\dots,c_9)=(12,\,132,\,324,\,498,\,504,\,336,\,144,\,36,\,4)\]
and $c_k=0$ for $k\ge10$. For $q=5$ the numbers are $5\cdot(12,132,325,498,504,336,144,36,4)$: the five
extra triples consist of the edges $\{B,C\}$ at the centres $a_i,a_{i+1},a_{i+3}$, which do not lie in three
consecutive units of the $5$-ring.
\end{theorem}

In particular a pair of edges is removable for only $132q$ of the $6q(12q-1)$ pairs, i.e.\ each edge of $L(J_q)$ can
be removed together with only $22$ others on average.

\section{The lower bound: claw sets}
\begin{proposition}\label{prop:claw}
Let $G$ be a snark and $v$ a vertex of $G$. Then the claw set $R_v$ is removable from $L(G)$.
\end{proposition}
\begin{proof}
Removing $R_v$ isolates the three vertices of $L(G)$ that are the edges at $v$, and what is left of $L(G)$ is
$L(G-v)$ plus three isolated vertices. So it suffices to show that $G-v$ has no $3$-edge-colouring. Suppose it
had one, and let $x,y,z$ be the neighbours of $v$; they have degree $2$ in $G-v$ and every other vertex has
degree~$3$. Each colour class is a matching, so the number of vertices it misses has the parity of
$|V(G-v)|$, which is odd because cubic graphs have an even number of vertices. The vertices missed by a colour
are among $x,y,z$, and each of them misses exactly one colour. Hence each of the three colours is missed by
exactly one of $x,y,z$, and colouring $vx,vy,vz$ with the colours missing at $x,y,z$ gives a $3$-edge-colouring
of $G$, a contradiction.
\end{proof}

The fifteen-edge set of Theorem~\ref{thm:main} for $q=3$ is removable for a similar reason. Removing it leaves
$L(J_3-\{b_0,b_1,b_2\})$ plus six isolated vertices. Contracting the triangle $b_0b_1b_2$ of $J_3$ gives the
Petersen graph $P$, so $J_3-\{b_0,b_1,b_2\}$ is $P$ minus a vertex, which has no $3$-edge-colouring by the
argument above applied to the snark $P$. More generally, if a snark $G$ has a triangle whose contraction gives a
simple graph $G'$, then $G'$ is again a snark, $G$ minus the triangle is $G'$ minus a vertex, and the fifteen
edges of $L(G)$ that touch the six edges at the triangle form a removable set.

\section{The upper bound: a transfer matrix}
\subsection*{Units and states}
Unit $i$ of $J_q$ consists of the vertices $a_i,b_i,c_i,d_i$, the spokes $A=a_ib_i$, $B=a_ic_i$, $C=a_id_i$, and the
three \emph{right connectors} $X=b_ib_{i+1}$, $Y=c_ic_{i+1}$, $Z=d_id_{i+1}$, except that for $i=q-1$ they are
$X=b_{q-1}b_0$, $Y=c_{q-1}d_0$, $Z=d_{q-1}c_0$. The right connectors of unit $i-1$ are the \emph{left connectors}
$X_L,Y_L,Z_L$ of unit $i$; at unit $0$ the roles of $Y$ and $Z$ are exchanged (the twist). A \emph{state} is a
triple of colours $(X,Y,Z)\in\{0,1,2\}^3$; there are $27$ states.

The twelve edges of $L(J_q)$ at the four vertices of unit $i$ get local numbers:
\begin{center}\small
\begin{tabular}{llll}
\toprule
at $a_i$ & at $b_i$ & at $c_i$ & at $d_i$\\
\midrule
$0:\{A,B\}$ & $3:\{X_L,A\}$ & $6:\{Y_L,B\}$ & $9:\{Z_L,C\}$\\
$1:\{A,C\}$ & $4:\{A,X\}$ & $7:\{B,Y\}$ & $10:\{C,Z\}$\\
$2:\{B,C\}$ & $5:\{X_L,X\}$ & $8:\{Y_L,Y\}$ & $11:\{Z_L,Z\}$\\
\bottomrule
\end{tabular}
\end{center}
For $S\subseteq\{0,\dots,11\}$ let $U(S)$ be the $27\times27$ Boolean matrix whose entry $(s,t)$ is $1$ when
there are colours for $A,B,C$ such that, with left connectors coloured by $s$ and right connectors by $t$, any two
edges of $J_q$ that form one of the twelve pairs get different colours unless the pair is in $S$. Let
$T=U(\emptyset)$, and let $W$ be the permutation matrix of $(X,Y,Z)\mapsto(X,Z,Y)$. A set $R$ of edges of $L(J_q)$
is a set of pairs (unit $i$, local number), and $R_i$ denotes its part in unit~$i$.

\begin{lemma}\label{lem:trace}
$L(J_q)-R$ is $3$-colourable if and only if the Boolean matrix
\[W\,U(R_0)\,U(R_1)\cdots U(R_{q-1})\]
has a $1$ on its diagonal.
\end{lemma}
\begin{proof}
A $3$-colouring of $L(J_q)-R$ is a colouring of the edges of $J_q$ in which two edges at a common vertex get
different colours unless the corresponding edge of $L(J_q)$ is in~$R$. Every such pair lies in exactly one unit,
and the constraints of unit $i$ involve only its spokes and its left and right connectors. So colourings
correspond to closed walks $s_{q-1}\to s_0\to\dots\to s_{q-1}$ through the states of the right connectors, the
step into unit $0$ going through the twist~$W$.
\end{proof}

The next facts are finite computations with $27\times27$ Boolean matrices (\S\ref{sec:check}).

\begin{lemma}\label{lem:facts}
\textup{(a)} $T^2=T^4$, so $T^k$ for $k\ge2$ depends only on $k\bmod2$.

\noindent\textup{(b)} $WTW=T$, and $W\,U(S)\,W=U(\sigma S)$ for every $S$; here $\sigma$ exchanges $0$ and $1$,
$6$ and $9$, $7$ and $10$, $8$ and $11$, and fixes $2,3,4,5$.
\end{lemma}

The map that sends unit $i$ to unit $i+1$ for $i<q-1$, and unit $q-1$ to unit $0$ with $c$ and $d$ exchanged,
is an automorphism of $J_q$. In terms of matrices this is (b) together with the cyclic invariance of the trace:
the twist can be moved past a unit at the cost of relabelling that unit by~$\sigma$. So any configuration can be
rotated so that a chosen unit becomes unit~$0$; the rotated configuration is again a set of (unit, local number)
pairs, of the same size.

\begin{lemma}\label{lem:far}
Let $q\ge7$ be odd. If two edges of $L(J_q)$ lie in units at cyclic distance at least $3$, they do not form a
removable pair. (For $q=5$ no two units are at distance $3$.)
\end{lemma}
\begin{proof}
Rotate so that the first edge, $e$, is in unit $0$ and the second, $f$, in unit $\delta$ with $3\le\delta\le q-3$.
By Lemma~\ref{lem:trace} the pair is removable iff
$W\,U(\{e\})\,T^{\delta-1}\,U(\{f\})\,T^{q-1-\delta}$ has no $1$ on its diagonal. Both exponents are at least $2$ and
their sum $q-2$ is odd, so by Lemma~\ref{lem:facts}(a) the product takes only the two values with exponents
$(2,3)$ and $(3,2)$. For each of these and each of the $144$ choices of $e,f$ the diagonal contains a~$1$.
\end{proof}

\begin{proof}[Proof of Theorem~\ref{thm:main} for $q\ge7$]
Let $R$ be removable. Every two of its edges form a removable pair, so by Lemma~\ref{lem:far} their units are at
cyclic distance at most~$2$. These units lie in an arc of three consecutive units: rotate so that one of them is
unit $0$; all others are in $\{-2,\dots,2\}$, and a unit $x>0$ and a unit $y<0$ would be at cyclic distance
$\min(x-y,\,q-x+y)\ge3$ unless $x-y\le2$, because $q-x+y\ge q-4\ge3$. So the units of $R$ are in $\{-2,-1,0\}$,
$\{-1,0,1\}$ or $\{0,1,2\}$. Rotate them to units $0,1,2$. By Lemma~\ref{lem:trace} and Lemma~\ref{lem:facts}(a), $R$ is removable iff
\[W\,U(R_0)\,U(R_1)\,U(R_2)\,T^2\]
has no $1$ on its diagonal, because $q-3\ge4$ is even. This condition does not depend on $q$. We enumerate all
sets $R$ of the $36$ edges of units $0,1,2$ that satisfy it, by depth-first search that adds edges in a fixed order
and stops as soon as a set fails (subsets of removable sets are removable). The numbers of such sets of sizes
$0,1,\dots,9$ are
\[1,\ 36,\ 297,\ 678,\ 951,\ 882,\ 546,\ 222,\ 54,\ 6,\]
and there are none of size $10$. Of the six sets of size~$9$, four meet unit $0$; they are the claw sets of $a_0$,
$b_1$, $c_1$, $d_1$ (the program reads them off the explicit graph). Every removable set has a unique first unit,
the first unit of the shortest arc that contains its units (an arc of three units has a unique first unit
when $q\ge7$), so for each odd $q\ge7$ the removable nine-edge sets are the $4q$ claw sets
$R_{a_i},R_{b_{i+1}},R_{c_{i+1}},R_{d_{i+1}}$, $i\in\mathbb Z_q$, and there is no removable set with ten edges.
\end{proof}

\begin{proof}[Proof of Theorem~\ref{thm:main} for $q=3,5$]
Here the same depth-first search runs over all $12q$ edges of the ring, using Lemma~\ref{lem:trace} to test each
set. For $q=5$ the largest removable sets have nine edges and there are $20$ of them; by
Proposition~\ref{prop:claw} the $20$ claw sets are among them, so these are all. For $q=3$ the largest removable
set has fifteen edges and is unique.
\end{proof}

\begin{proof}[Proof of Theorem~\ref{thm:count}]
As in the proof above, a removable $k$-set of $L(J_q)$, $q\ge7$ odd, corresponds to a removable $k$-set of the
window of units $0,1,2$ that meets unit~$0$, rotated to each of the $q$ possible first units. The number of window
sets meeting unit $0$ is the number of removable window sets minus the number of those that avoid unit $0$.
The window sets avoiding unit $0$, i.e.\ the removable sets inside units $1,2$ alone, number
$1,24,165,354,453,378,210,78,18,2$ for $k=0,\dots,9$; subtracting these from the window counts gives the $c_k$.
The counts for $q=5$ come from the exhaustive search over the whole ring.
\end{proof}

\section{Checks}\label{sec:check}
The program \texttt{snark\_check.py} (Python standard library only) does the following.
\begin{itemize}
\item \texttt{model}: for $q=3,5,7$ and for every set $R$ of at most two edges of $L(J_q)$ ($6069$ sets in all) it
compares Lemma~\ref{lem:trace} with a direct backtracking $3$-colouring of the explicit graph $L(J_q)-R$; it also
checks that the local numbering is a bijection onto the $12q$ edges. No disagreement.
\item \texttt{period}: $T^2=T^4$ (and $T\ne T^3$); $WTW=T$; $W\,U(S)\,W=U(\sigma S)$ for all $4096$ sets $S$.
\item \texttt{pairs}: the computation in the proof of Lemma~\ref{lem:far}.
\item \texttt{window}: the enumeration in the proof of Theorem~\ref{thm:main}, including the identification of the
four maximum sets with claw sets of the explicit graph.
\item \texttt{small}: the cases $q=3,5$, including the description of the fifteen-edge set for $q=3$ and the
counts for $q=5$.
\item \texttt{claws}: every claw set is removable for all odd $q\le41$ by Lemma~\ref{lem:trace}, and for $q\le9$ also
by direct colouring.
\end{itemize}
Independently of the transfer matrix, a CP-SAT model with the removable pairs as constraints, followed by a direct
colourability test of each candidate, gives the maximum $9$ with exactly $4q$ maximum sets for $q=5,7,\dots,15$,
the maximum $15$ for $q=3$, and $132q$ removable pairs and $324q$ removable triples for $q=5,\dots,11$ and $q=7,9$
respectively, in agreement with Theorem~\ref{thm:count}. We also checked that our $J_q$ is isomorphic to the graph
of~\cite[\S7.2.2.2, ex.~176]{knuth4b} for $q=3,5,7,9$. Finally, an independent referee (another AI agent) rebuilt
the graphs, a SAT-based colourability test and the transfer matrices from the text of this note alone, and
found Theorems~\ref{thm:main} and~\ref{thm:count} confirmed by exhaustive enumeration for every odd
$q\le17$.

\section{Remarks}
For the Petersen graph $P$ the same search shows that the largest removable sets of $L(P)$ also have nine edges,
and that there are exactly ten of them, the claw sets of its ten vertices. By Proposition~\ref{prop:claw} nine
edges can be removed from $L(G)$ for every snark $G$, and fifteen when $G$ has a suitable triangle
(\S2). The graph $J_3$ shows that more than nine is sometimes possible.

\begin{question}
For which snarks $G$ is nine the largest number of removable edges of $L(G)$, with the claw sets as the only
nine-edge removable sets? Is this true for all snarks of girth at least five?
\end{question}

\medskip\noindent\emph{How this was done.} The search programs, the transfer-matrix argument, the checks and this
note were produced with Claude Code (Anthropic) at the author's direction. The author is responsible for the
content.

\begin{thebibliography}{9}
\bibitem{isaacs} R.~Isaacs, Infinite families of nontrivial trivalent graphs which are not Tait colorable,
\emph{Amer. Math. Monthly} \textbf{82} (1975), 221--239.
\bibitem{knuth4b} D.~E. Knuth, \emph{The Art of Computer Programming}, Volume 4B, Addison--Wesley, 2022.
\bibitem{knuth7} D.~E. Knuth, \emph{The Art of Computer Programming}, Volume 4, Fascicle 7: Constraint
Satisfaction, Addison--Wesley, 2025. Cited from its final draft, Pre-Fascicle 7A (5 December 2024), exercise
7.2.2.3--233 and its answer.
\end{thebibliography}
\end{document}
